% BONFERRONNIMOMENT1
\documentclass[11pt]{article}
\usepackage[margin=.88in]{geometry}
\usepackage{amsmath,amssymb,amsthm}
\newtheorem{theorem}{Theorem}
\newtheorem{corollary}{Corollary}
\begin{document}

\begin{center}
{\Large BONFERRONNI-MOMENT1: Exact DNF fooling from slice moment matching}\\[4pt]
\textbf{Date: 2026-10-04 22:30 PT}\\
\textbf{Status: PROVED/EXACT lemmas + COMPUTATIONAL hostile control}\\
\textbf{P versus NP status: OPEN}
\end{center}

\section{Literature position}

Wegener's slice theory identifies set circuits as the natural exact-slice
model and warns that the standard monotone lower-bound methods do not transfer
to slice functions in their original form.

Cavalar--Kumar--Rossman (Algorithmica 2022) prove lower bounds for general
monotone DAG circuits by replacing each gate with a closed-and-trimmed
approximator. Their analysis controls positive errors through small minterms
and negative errors through closure.  The paper also states that the
classical sunflower approximation framework leaves limited room for further
improvement when used with the same product-style test distributions.

Beimel--Nir (ECCC TR26-213, 2026) obtain strong lower bounds for
\emph{formulas} over multislice gates.  This is not an unrestricted monotone
DAG lower bound and therefore does not close TASK-008.

The more recent method-of-moments sunflower theorem (ECCC TR26-220, 2026) is
useful conceptually, but SLICEWIDTH1 and HYPERGEOMCENTER1 show that its
published pseudo-independence hypothesis and centered-factorization route do
not directly fit the present exact sparse slice.

This note extracts the method-of-moments step that \emph{does} transfer
exactly.

\section{Exact Bonferroni moment lemma}

Let \(\mu_+\) and \(\mu_-\) be distributions on
\(\{0,1\}^M\).  Assume they match every squarefree coordinate moment through
degree \(K\):
\[
\mathbb E_{\mu_+}\prod_{i\in S}z_i
=
\mathbb E_{\mu_-}\prod_{i\in S}z_i
\qquad
\forall |S|\le K.
\tag{1}
\]

Let
\[
H(z)=\bigvee_{a=1}^m \bigwedge_{i\in A_a} z_i
\]
be a monotone DNF, with \(|A_a|\le w\).  Define
\[
X(z)=\sum_{a=1}^m
\mathbf 1[A_a\subseteq z].
\]

\begin{theorem}[BONFERRONNI-MOMENT]
Let \(R\ge2\) be even and satisfy \(Rw\le K\).  Then
\[
\boxed{
\left|
\Pr_{\mu_+}[H=1]
-
\Pr_{\mu_-}[H=1]
\right|
\le
\mathbb E_{\mu_-}\binom XR
}
\tag{2}
\]
and the same factorial moment occurs under \(\mu_+\).
\end{theorem}

\begin{proof}
For every \(j\le R\),
\[
\mathbb E\binom Xj
=
\sum_{I\subseteq[m], |I|=j}
\Pr\left[
\bigcup_{a\in I}A_a\subseteq z
\right].
\]
Every union has size at most \(jw\le Rw\le K\), so (1) implies equality
of all factorial moments through order \(R\).

Bonferroni gives, for even \(R\),
\[
\sum_{j=1}^{R}(-1)^{j+1}\mathbb E\binom Xj
\le
\Pr[X>0]
\le
\sum_{j=1}^{R-1}(-1)^{j+1}\mathbb E\binom Xj.
\]
The two distributions therefore lie in the same interval, whose width is
\(\mathbb E\binom XR\).  This proves (2).
\end{proof}

This is an exact two-distribution statement and uses only raw containment
moments.  It is unaffected by CENTER-RATIO-FAIL.

\section{Uniform exact-slice factorial moment}

For
\[
\mu_-=U_{M,t},
\]
the uniform exact-\(t\) slice,
\[
\mathbb E_{U_{M,t}}\binom XR
=
\sum_{I\subseteq[m], |I|=R}
\frac{(t)_{|\cup_{a\in I}A_a|}}
{(M)_{|\cup_{a\in I}A_a|}}.
\tag{3}
\]

Thus the entire DNF fooling error is an explicit combinatorial quantity.

\section{Literal sunflower corollary}

Suppose the DNF terms form a literal sunflower
\[
A_a=C\cup P_a,
\qquad
|C|=c,
\qquad
|P_a|=d,
\]
with pairwise disjoint petals.  Then every \(R\)-term union has size
\[
c+Rd.
\]

\begin{corollary}[SUNFLOWER-MOMENT]
If \(R\) is even and \(c+Rd\le K\), then
\[
\boxed{
\left|
\Pr_{\mu_+}[H=1]
-
\Pr_{U_{M,t}}[H=1]
\right|
\le
\binom mR
\frac{(t)_{c+Rd}}{(M)_{c+Rd}}
}.
\tag{4}
\]
\end{corollary}

This removes the term-by-term union bound that caused DENSITY-CANCEL1.
The required resource is now growing moment order, not a product-measure
centered-factorization theorem.

\section{Numerical parameter checks}

At density \(t=3.5q\):

\[
\begin{array}{c|c|c|c|c|c}
q&c&d&m&R&\text{bound in (4)}\\
\hline
20&0&3&300000&4&5.23\times10^{-6}\\
24&0&3&300000&4&5.65\times10^{-8}\\
20&1&2&10000&4&2.26\times10^{-5}\\
24&1&2&10000&4&7.29\times10^{-7}\\
20&2&1&1000&10&4.08\times10^{-3}\\
24&2&1&1000&10&4.41\times10^{-5}
\end{array}
\]

The corresponding moment orders \(K=c+Rd\) are 12, 12, 9, 9, 12, and 12.

These are exact-slice combinatorial bounds for the stated literal sunflower
DNFs, not empirical acceptance estimates.

\section{XOR-SLICE-EASY hostile control}

Moment camouflage itself is not SAT-specific.

At density \(0.9\), where both outcomes occur for the Gaussian-elimination
parity-shadow predicate, exact canonical three-clause orbit vectors were
computed for every sampled formula.

For \(q=12,t=11,N=5000\):
\[
\#XOR\text{-UNSAT}=1728,
\qquad
\#XOR\text{-SAT}=3272,
\]
and the empirical XOR-SAT target was at numerical \(L_\infty\) distance
\[
3.70\times10^{-6}
\]
from the convex hull of XOR-UNSAT order-three vectors.

A second \(q=12,N=10000\) run gave residual
\[
2.33\times10^{-6}.
\]

At \(q=14,t=13,N=5000\):
\[
\#XOR\text{-UNSAT}=1847,
\qquad
\#XOR\text{-SAT}=3153,
\]
and the numerical LP returned
\[
\boxed{\delta=0}
\]
to solver tolerance, with 68 positive support columns.

Because the parity-shadow predicate is polynomial-time solvable by Gaussian
elimination, low-order fractional-design feasibility cannot by itself be a
hardness witness.

\section{What this changes}

The missing bridge should no longer be phrased as
\[
\text{small circuit}\Rightarrow\text{low degree}.
\]

The more precise target is:

\[
\boxed{
\text{small UNSAT set circuit}
\Rightarrow
\text{two-distribution approximator decomposable into bounded-width
sunflower DNFs with small factorial tails}
}.
\tag{5}
\]

If (5) were proved in a SAT-specific way, then growing-order UNSAT fractional
designs plus BONFERRONNI-MOMENT would control the approximator without using
the pseudo-independent product-mixture theorem.

The SAT-specific qualifier is essential: XOR-SLICE-EASY shows that the local
moment side of the argument can also hold for an easy predicate.

Krieger's prime-implicant incompressibility results are relevant here because
they exploit prime-implicant structure and apply to bounded-depth/multilinear
models and pseudoslice-like functions.  They do not provide (5) for general
monotone DAG circuits.

\section{Next target}

\textbf{SETCIRCUIT-SUNFLOWER1.}

Using Wegener's set-circuit representation of a \(t\)-slice, formulate the
output gate directly as a family of exact-layer positive sets.  Attempt to
prove or falsify a gate-by-gate decomposition theorem in which every
intermediate family either

\begin{enumerate}
\item contains a large literal/robust sunflower whose acceptance is controlled
by BONFERRONNI-MOMENT, or
\item has a quantitatively small family of bounded-width minterms and can be
charged directly.
\end{enumerate}

The theorem must be strong enough for DAG sharing, not just formulas, and it
must fail cleanly on XOR-SLICE-EASY unless it uses SAT-specific core/minterm
geometry.

In parallel, continue DUAL-COLUMN3 only as a way to test growing moment order;
order-three feasibility alone is now known not to be SAT-specific.

\section{Status}

BONFERRONNI-MOMENT: \textbf{PROVED/EXACT}.

SUNFLOWER-MOMENT: \textbf{PROVED/EXACT}.

XOR order-three fractional-design control: \textbf{COMPUTATIONAL}.

SETCIRCUIT-SUNFLOWER1: \textbf{OPEN}.

General monotone slice lower bound and \(P\ne NP\): \textbf{OPEN}.

\begin{thebibliography}{5}
\bibitem{Wegener1985}
I. Wegener,
On the Complexity of Slice Functions,
Theoretical Computer Science 38:55--68, 1985.

\bibitem{CKR2022}
B. P. Cavalar, M. Kumar, and B. Rossman,
Monotone Circuit Lower Bounds from Robust Sunflowers,
Algorithmica 84:3655--3685, 2022.

\bibitem{CFMSY2026}
B. P. Cavalar, T. Fabris, P. Mukhopadhyay, S. Srinivasan, and A. Yehudayoff,
Derandomized Sunflowers and Radical Lower Bounds,
ECCC TR26-220, 2026.

\bibitem{BeimelNir2026}
A. Beimel and O. Nir,
The Power of Multislices in Monotone Computation,
ECCC TR26-213, 2026.

\bibitem{Krieger2007}
M. P. Krieger,
On the Incompressibility of Monotone DNFs,
Theory of Computing Systems 41:211--231, 2007.
\end{thebibliography}

\end{document}
